\documentclass[aspectratio=169,11pt]{beamer}
\usetheme{Madrid}
\usecolortheme{beaver}
\setbeamertemplate{navigation symbols}{}

\usepackage{listings}
\usepackage{booktabs}
\usepackage{xcolor}

\definecolor{codebg}{RGB}{245,245,245}
\definecolor{good}{RGB}{0,120,60}
\definecolor{bad}{RGB}{180,0,0}
\lstset{
  language=Python,
  basicstyle=\ttfamily\small,
  backgroundcolor=\color{codebg},
  keywordstyle=\color{blue!70!black}\bfseries,
  commentstyle=\color{gray}\itshape,
  stringstyle=\color{orange!70!black},
  frame=single, rulecolor=\color{gray!40},
  showstringspaces=false, columns=fullflexible,
  aboveskip=4pt, belowskip=4pt
}

\title[ADV TP1 Summary]{ADV TP1: What I Understood}
\subtitle{Advanced Programming for Data Science}
\author[HEM Bellyday]{HEM Bellyday}
\institute[ITC]{Department of Applied Mathematics and Statistics (AMS)\\
Institute of Technology of Cambodia (ITC)\\[2pt]
Lecturer: Mr.~VANN Visal}
\date{Lecture of October 13, 2026}

\begin{document}

\begin{frame}
  \titlepage
\end{frame}

\begin{frame}{Outline}
  \tableofcontents
\end{frame}

% ===================== PART 1 =====================
\section{Part 1: Building Blocks}

\begin{frame}{Part 1: Building Blocks (1/2)}
  \begin{block}{1. What is a function?}
    This is like a math function: it has something like $x$ or $y$ in it when we
    define it. We use it so we don't do a task again and again; we just use this
    function and finish. (Write once, use quickly.)
  \end{block}
  \begin{block}{2. What is a variable?}
    A variable is like a math definition: it refers to a value or object when we
    want to define it.
  \end{block}
  \begin{block}{3. What is a class?}
    A class is like a function, but it has a bunch of functions inside of it.
  \end{block}
\end{frame}

\begin{frame}{Part 1: Building Blocks (2/2)}
  \begin{block}{4. What is a method?}
    A method is a function that makes something work, or does an action, inside a
    class. Example on the board: calculate score.
  \end{block}
  \begin{block}{5. What is a module?}
    A module is like a \texttt{.py} file: we write functions inside it and call
    them to use in the future in another module.
  \end{block}
  \begin{block}{6. What is a package?}
    A bunch of modules stuck together. When we install a package, that is a bunch
    of modules ready to use.
  \end{block}
  \begin{block}{7. What is a constant?}
    It refers to a value that doesn't change, even in future use.
  \end{block}
\end{frame}

% ===================== PART 2 =====================
\section{Part 2: Clean Code}

\begin{frame}{Part 2: Clean Code}
  \begin{alertblock}{Before clean code}
    The teacher asked what is something in class we don't understand, and I said
    PEP~8.
  \end{alertblock}
  \begin{block}{What is clean code?}
    Clean code refers to a file or code lines written so that not only the user
    can understand them, but others can too.
  \end{block}
\end{frame}

\begin{frame}{Part 2: Clean Code Practices}
  \begin{itemize}\setlength\itemsep{6pt}
    \item[10.] \textbf{Descriptive name:} name the variable so it can be
      understood quickly, and keep it short. Example: a function to start a moto
      engine; there are many names we could use, but use one name that other
      functions in the moto can't use.
    \item[11.] \textbf{Comment:} \texttt{\#} that explains code in a single line.
    \item[12.] \textbf{Docstring:} used to explain a whole function or class so
      others can understand it.
    \item[13.] \textbf{Naming values:} name values that can't be changed (even in
      the future they don't change), so it is easy and we don't have to find that
      value in whatever file it is in.
    \item[14.] \textbf{Functions:} a function needs to be clear and reasonable to
      use.
  \end{itemize}
\end{frame}

% ===================== PART 3 =====================
\section{Part 3: PEP 8}

\begin{frame}{Part 3: What is PEP 8?}
  \begin{block}{PEP 8}
    PEP~8 is created by Python itself, just to make code structure more clean and
    reasonable, and it is also a rule to follow from Python.
  \end{block}
  \vspace{6pt}
  \centering
  \begin{tabular}{@{}lll@{}}
    \toprule
    \textbf{What} & \textbf{Style} & \textbf{Example} \\
    \midrule
    Variables, functions, methods & \texttt{snake\_case} & \texttt{student\_score} \\
    Constants & \texttt{UPPER\_CASE} & \texttt{PASS\_MARK} \\
    Classes & \texttt{PascalCase} & \texttt{Student} \\
    \bottomrule
  \end{tabular}
\end{frame}

\begin{frame}{Part 3: PEP 8 Rules in Simple Words}
  \small
  \begin{itemize}\setlength\itemsep{3pt}
    \item \textbf{Indentation:} 4 spaces for each level.
    \item \textbf{Line length:} at most 79 characters for code, 72 for comments
      and docstrings.
    \item \textbf{Imports:} put them at the top. First Python's built-in modules,
      then installed libraries, then my own files, with a blank line between each
      group. Don't use \texttt{from x import *}.
    \item \textbf{Spaces:} one space around \texttt{=} and \texttt{==}, but no
      extra space just inside brackets or before the \texttt{(} of a call.
    \item \textbf{Commas and colons:} a space after a comma, like
      \texttt{[60, 80]}, but no spaces in a slice, like \texttt{scores[1:3]}.
    \item \textbf{Blank lines:} 2 between top-level functions or classes, 1
      between methods inside a class.
    \item \textbf{Long lines:} wrap them inside parentheses, and break the line
      before an operator like \texttt{+} or \texttt{and}.
    \item \textbf{Quotes:} choose \texttt{'} or \texttt{"} and keep using the same
      one. Docstrings use \texttt{"""}.
  \end{itemize}
\end{frame}

\begin{frame}[fragile]{Part 3: PEP 8 Example}
  \begin{columns}[T]
    \begin{column}{0.48\textwidth}
      \textcolor{bad}{\textbf{Not following PEP 8}}
\begin{lstlisting}
passmark=50
def HasPassed( s ):
  return s>=passmark
print (HasPassed(75))
\end{lstlisting}
    \end{column}
    \begin{column}{0.48\textwidth}
      \textcolor{good}{\textbf{Following PEP 8}}
\begin{lstlisting}
PASS_MARK = 50


def has_passed(score):
    return score >= PASS_MARK


print(has_passed(75))
\end{lstlisting}
    \end{column}
  \end{columns}
\end{frame}

% ===================== PART 4 =====================
\section{Part 4: Debugging}

\begin{frame}{Part 4: Debugging}
  \begin{block}{What is debugging?}
    Fixing errors, of course. But sometimes the error shown is not clear or
    reasonable, so we have to check further in the code file.
  \end{block}
  \begin{block}{Tools the teacher recommended}
    The teacher recommended 2 old-school tools. Even before AI existed, these are
    what old programmers used. (I forgot the names.)
  \end{block}
\end{frame}

\begin{frame}{Part 4: How I Debug}
  \begin{enumerate}\setlength\itemsep{8pt}
    \item Print the suspicious code line, so we can clearly see the error and know
      what to fix.
    \item Don't put 5 million lines of code in 1 file. Divide them and use them as
      imports. When there is an error, it is easy to find in the file structure.
  \end{enumerate}
  \vspace{16pt}
  \begin{center}
    \large\textit{That's all I learned from that day.}
  \end{center}
\end{frame}

\end{document}
